\documentclass[11pt,a4paper]{article}
\usepackage[utf8]{inputenc}
\usepackage{geometry}
\geometry{margin=1in}
\usepackage{hyperref}
\usepackage{enumitem}
\usepackage{amsmath}

\title{Digital Design Training Program \\ \large Day 7: Custom LCD Drivers and 8-Bit Architecture Limits}
\author{Author,who}
\date{April 1, 2026}

\begin{document}

\maketitle

\section*{Theme: Building a bare-metal communication driver for external displays and understanding the memory/processing constraints of an 8-bit microcontroller without a Floating Point Unit (FPU).}

\subsection*{Module 7.1: The 4-Bit Bare-Metal LCD Driver (120 Minutes)}

\begin{itemize}
    \item \textbf{Goal:} Write a complete C driver from scratch to command a 16x2 HD44780 LCD, mastering the concept of bit-shifting data into ``nibbles'' to save microcontroller pins.
    \item \textbf{Hardware Setup (The PORTB Remap):}
    \begin{itemize}
        \item To ensure compatibility with our final Capstone, the LCD \emph{must} avoid \verb|PORTD| and \verb|PORTC| (used by the Keypad). Wire the LCD to \verb|PORTB| (Digital Pins 8-13):
        \item \textbf{RS (Register Select):} \verb|PB0| (Pin 8)
        \item \textbf{EN (Enable):} \verb|PB1| (Pin 9)
        \item \textbf{D4 to D7 (Data Lines):} \verb|PB2| to \verb|PB5| (Pins 10 to 13)
        \item \textbf{R/W:} Tie directly to Ground (we only write to the LCD).
    \end{itemize}
    \item \textbf{Action Steps:}
    \begin{enumerate}[label=\arabic*.]
    \item \textbf{The Enable Pulse:} The LCD only reads data when the EN pin drops from HIGH to LOW. Write a \verb|lcd_pulse()| function. \emph{Note: The HD44780 requires strict $\sim$1 microsecond delays to latch data. For this specific low-level hardware timing, we will allow the use of} \verb|<util/delay.h>| (\verb|_delay_us(1)|).
        \item \textbf{Nibble Shifting:} Because we are using 4-bit mode (saving 4 wires), an 8-bit character (like `A', \verb|0x41|) must be sent in two halves. Write a function that takes an 8-bit \verb|char|, shifts the top 4 bits into \verb|PB2-PB5|, pulses EN, then shifts the bottom 4 bits into \verb|PB2-PB5|, and pulses EN again.
        \item \textbf{Initialization Sequence:} The hardest part of custom drivers is the ``wake-up'' sequence. Follow the datasheet to send the specific hex commands (\verb|0x33|, \verb|0x32|, \verb|0x28|, \verb|0x0C|, \verb|0x01|) using \verb|lcd_command()|.
        \item \textbf{Printing:} Write an \verb|lcd_print(char* str)| function that loops through a character array and passes each letter to \verb|lcd_data()| until it hits the null terminator (\verb|'\0'|).
    \end{enumerate}
    \item \textbf{Checkpoint:} The screen successfully displays ``Hello World!'' purely through bare-metal register manipulation, without \verb|#include <LiquidCrystal.h>|.
\end{itemize}
\newpage
\subsection*{Module 7.2: 8-Bit Math Limits \& PROGMEM (30 Minutes)}

\begin{itemize}
    \item \textbf{Goal:} Understand the cost of software-emulated floating-point math on an AVR chip, and learn how to utilize Flash memory to save RAM.
    \item \textbf{Action Steps:}
    \begin{enumerate}[label=\arabic*.]
        \item \textbf{The Emulation Cost:}The ATmega328P has no hardware to do decimal math. When you type \verb|float z = x * y;|, the compiler secretly injects hundreds of lines of assembly code to emulate it. To see this in action, compile a blank script, then add a float multiplication and compile again to observe the massive jump in binary size.
        \item \textbf{SRAM vs. Flash:} The Arduino Uno only has 2KB of SRAM (variables) but 32KB of Flash (program space).
        \item \textbf{PROGMEM Implementation:} We introduce \verb|<avr/pgmspace.h>|. Write a large array of constant \verb|float| values (e.g., standard conversion factors) and use the \verb|PROGMEM| macro to force the compiler to store it in Flash memory. Read it back using \verb|pgm_read_float()|.
    \end{enumerate}
    \item \textbf{Checkpoint:} Understanding how to offload read-only data from precious SRAM and to be economical with \verb|float| types.
\end{itemize}
\newpage
\subsection*{Module 7.3: Mathematical Modeling via ODEs (45 Minutes)}

\begin{itemize}
    \item \textbf{Goal:} Shift the focus from hardware to numerical methods. Learn how to reframe complex functions (sine, exponential, logarithms) into initial value problems that a computer can solve through repeated addition.
    \item \textbf{Action Steps:}
    \begin{enumerate}[label=\arabic*.]
        \item \textbf{The Math Void:} We cannot easily use standard \verb|<math.h>| libraries for our scientific calculator because they consume too much memory and processing time. We must build our own math engine.
        \item \textbf{Defining the ODEs:} Frame target functions as Ordinary Differential Equations (ODEs) with an Initial Value:
        \begin{itemize}
            \item To find $y=e^{x}$, we solve the ODE $y^{\prime}=y$, where $y(0)=1$.
            \item To find $y=\sin(x)$, we solve a system where $y^{\prime}=\cos(x)$ and $y(0)=0$.
        \end{itemize}
    \end{enumerate}
    \item \textbf{Checkpoint:} Understanding the concept of numeric integration.
\end{itemize}

\end{document}
